\chapter{Giới thiệu}
\label{chap:introduction}

\section{Mục đích}

Tài liệu mô tả thiết kế giải pháp SmartWater theo trạng thái của repository tại thời
điểm khảo sát. Nội dung cung cấp kiến trúc, ranh giới thành phần, luồng telemetry,
thiết kế dữ liệu, điểm tích hợp, bảo mật, vận hành, trạng thái triển khai và các khoảng
trống có bằng chứng. Tài liệu không thay thế kiểm thử tích hợp hoặc xác nhận môi trường
vận hành.

\section{Phạm vi}

Phạm vi bao gồm firmware ESP32 và simulator, cấu hình client HiveMQ/MQTT, Device
Monitor, SmartWater Admin, migration MySQL dùng chung, worker .NET và các giao diện
liên quan. Broker HiveMQ và máy chủ MySQL được xem là phụ thuộc bên ngoài khi repository
không cung cấp manifest triển khai tương ứng.

\section{Nguồn bằng chứng và nguyên tắc}

Các nguyên tắc được sử dụng xuyên suốt tài liệu như sau:

\begin{itemize}
    \item \textbf{Migration là nguồn sự thật (Source of Truth).}
    Khi có khác biệt giữa Model, DDL cục bộ hoặc câu lệnh SQL, các migration trong thư mục
    \texttt{database/migrations} được xem là nguồn tham chiếu chính thức.

    \item \textbf{Khóa định danh telemetry.}
    Khóa định danh chuẩn (canonical identifier) của dữ liệu telemetry là
    \texttt{mcu\_id}.

    \item \textbf{Xử lý khác biệt lịch sử.}
    Các thành phần vẫn sử dụng \texttt{device\_id}, schema
    \texttt{device\_data} hoặc MQTT topic cũ được xem là ranh giới tích hợp
    (integration boundary), không tự động hợp nhất hoặc diễn giải bằng suy đoán.

    \item \textbf{Đối chiếu bằng chứng.}
    Danh mục bằng chứng và vị trí mã nguồn tương ứng được trình bày tại
    Phụ lục~\ref{app:source-evidence}.
\end{itemize}
\section{Đối tượng sử dụng tài liệu}

Tài liệu phục vụ kiến trúc sư giải pháp, nhóm phát triển firmware/backend/frontend,
nhóm kiểm thử, quản trị cơ sở dữ liệu, vận hành hệ thống và người phê duyệt kỹ thuật.
Người đọc cần dùng trạng thái triển khai kèm evidence path để phân biệt source đã có
với tích hợp đã được xác minh runtime.

\section{Thuật ngữ và chữ viết tắt}

\begin{longtable}{p{0.18\textwidth}p{0.72\textwidth}}
\toprule
\textbf{Thuật ngữ} & \textbf{Ý nghĩa trong tài liệu} \\
\midrule
MCU & Bộ điều khiển vi xử lý; \texttt{mcu\_id} là định danh canonical dạng chuỗi của telemetry. \\
TDS & Total Dissolved Solids, đại lượng chất rắn hòa tan được lưu trong telemetry. \\
MQTT & Giao thức publish/subscribe giữa publisher, broker và consumer. \\
WSS & WebSocket Secure dùng bởi mqtt.js trong Device Monitor. \\
API & Giao diện HTTP/JSON của Device Monitor hoặc SmartWater Admin. \\
ERD & Sơ đồ quan hệ thực thể được dựng theo migration canonical. \\
Canonical & Hợp đồng hoặc schema được chọn làm nguồn chuẩn hiện tại. \\
Legacy & Hợp đồng \texttt{device\_id}/\texttt{device\_data} chưa tương thích canonical. \\
\bottomrule
\end{longtable}

\section{Tài liệu tham chiếu}

\begin{itemize}
  \item Migration canonical: \path|Project/smartwater-database/database/migrations|.
  \item Tài liệu database và telemetry: \path|Project/smartwater-database/README.md|.
  \item Manifest firmware/simulator: hai file \path|platformio.ini| tương ứng.
  \item Manifest Device Monitor và Admin: \path|composer.json|, \path|package.json|.
  \item Project worker .NET: \path|Hivemq-Service/SmartWater.MqttService/SmartWater.MqttService.csproj|.
  \item Chỉ mục evidence đầy đủ: Phụ lục~\ref{app:source-evidence}.
\end{itemize}

\section{Cấu trúc tài liệu}

Các Chương~\ref{chap:solution-overview}--\ref{chap:data-design} trình bày tổng quan,
kiến trúc, thành phần và dữ liệu. Các chương tiếp theo phân tích tích hợp, Admin, bảo
mật, vận hành, công nghệ, triển khai và khả năng mở rộng. Phần cuối tổng hợp trạng thái,
rủi ro, roadmap và kết luận. Các phụ lục cung cấp inventory, catalogue, đặc tả hợp đồng,
bằng chứng hình ảnh và ma trận chức năng.
